% Academic talk: overlap-dependent CG iteration growth in ASM(BASIC)+ICC.
% Minimal style, complete tables.  Compile with XeLaTeX (ctex/Fandol).
\PassOptionsToPackage{table}{xcolor}
\documentclass[aspectratio=169,10pt]{beamer}

\usepackage[fontset=fandol]{ctex}
\usepackage{amsmath,amssymb}
\usepackage{booktabs}
\usepackage{array}

\definecolor{Struct}{HTML}{1F3A5F}   % single restrained structure colour
\definecolor{Gray}{HTML}{6E6E6E}

\usetheme{default}
\usefonttheme{professionalfonts}
\setbeamertemplate{navigation symbols}{}
\setbeamercolor{structure}{fg=Struct}
\setbeamercolor{frametitle}{fg=Struct}
\setbeamerfont{frametitle}{series=\bfseries,size=\large}
\setbeamercolor{normal text}{fg=black}
\setbeamertemplate{itemize item}{\textbullet}
\setbeamertemplate{itemize subitem}{--}
\setbeamertemplate{blocks}[default]
\setbeamercolor{block title}{fg=white,bg=Struct}
\setbeamercolor{block body}{bg=Struct!6}
% minimal frame title (no decorative rule), plain footer with page number
\setbeamertemplate{frametitle}{\vskip0.5em\hspace{0.2em}\usebeamerfont{frametitle}\insertframetitle%
  {\ifx\insertframesubtitle\empty\else\\[0.1em]\small\color{Gray}\insertframesubtitle\fi}\par}
\setbeamertemplate{footline}{\hfill\footnotesize\color{Gray}\insertframenumber\,/\,\inserttotalframenumber\hspace{0.4cm}\vspace{3pt}}

\newcommand{\ua}{$\uparrow$}
\newcommand{\da}{$\downarrow$}
\renewcommand{\arraystretch}{1.15}

\title{重叠型加性 Schwarz(BASIC)+ICC 预条件下\\ CG 迭代次数随 overlap 增大而增加的现象}
\subtitle{现象的复现、子域求解精度与重叠加权的作用,及单层修复}
\date{}
\author{}

\begin{document}

% ---- 1. Title ----
\begin{frame}[plain]
  \vskip1.2cm
  {\color{Struct}\large\bfseries 重叠型加性 Schwarz(BASIC)$+$ICC 预条件下\\[0.3em]
   CG 迭代次数随 overlap 增大而增加的现象\par}
  \vskip0.6em
  {\color{Gray} 现象的复现、子域求解精度与重叠加权的作用,及单层修复\par}
  \vskip1.0em
  \hrule height0.4pt
  \vskip0.8em
  \footnotesize
  模型问题:立方体 $[0,1]^3$ 上的 P1 有限元椭圆方程\\[0.3em]
  实现:MFEM 4.9 + PETSc 3.24,以及独立的纯 PETSc 实现(交叉验证)\\[0.3em]
  代码与数据:\texttt{github.com/tianhaomahpc-ops/cg-asm-icc-research}
\end{frame}

% ---- 2. Background / observation ----
\begin{frame}{背景与观察到的现象}
\textbf{背景.} 在心脏电生理代码(cardioid)中,椭圆型子系统用
\[
\text{CG}\ +\ \text{加性 Schwarz(PCASM,BASIC 变体)}\ +\ \text{子域不完全 Cholesky(ICC)}
\]
求解。增大子域 overlap 时,观察到 CG 迭代次数不降反升。\par
\vskip0.8em
\textbf{观察(原始数据,torso 子系统,ICC 填充 $L{=}0$).}
\vskip0.4em
\begin{center}
\begin{tabular}{ccc}
\toprule
overlap & 平均迭代次数 & 收敛 \\
\midrule
$O=0$ & 115 & 收敛 \\
$O=1$ & 200 & 达到迭代上限 \\
$O=2$ & 200 & 达到迭代上限 \\
\bottomrule
\end{tabular}
\end{center}
\vskip0.6em
经典重叠 Schwarz 理论(子域精确求解)预期:overlap 增大使条件数减小,迭代次数\emph{减少}。
此处与预期相反。本报告在干净的模型问题上复现该现象,辨析其来源,并给出单层修复。
\end{frame}

% ---- 3. Model problem & discretisation ----
\begin{frame}{模型问题与离散}
\textbf{主测问题(贯穿全部结果):椭圆方程,混合边界.}
\[
-\Delta u = 1\ \ \text{于}\ \Omega=[0,1]^3,\qquad
u=0\ \ \text{于}\ x{=}0\ \text{面},\qquad
\partial u/\partial n=0\ \ \text{于其余 5 面}.
\]
\vskip0.4em
\begin{itemize}\itemsep0.35em
  \item \textbf{网格.} 结构化 $n_x\times n_x\times n_x$ 六面体;每个六面体按 MFEM 的
        \texttt{hex\_to\_tet} 分解(共用主对角线)剖为 6 个四面体。
  \item \textbf{单元.} P1(线性)$H^1$ 元;自由度位于顶点,共 $(n_x{+}1)^3$ 个。
  \item \textbf{规模.} 结果主要取 $n_x=48$($117\,649$ 自由度),4 个 MPI 进程,METIS 图剖分。
  \item \textbf{解析参考解.} 由对称性 $u=x-\tfrac12 x^2$,用于校验数值解的正确性。
\end{itemize}
\vskip0.5em
另取两个问题作普遍性检验:全 Neumann(奇异)椭圆,以及含质量项的
reaction--diffusion 矩阵 $\tfrac1{\Delta t}M+\tfrac12 K$。
\end{frame}

% ---- 4. Solver & parameters ----
\begin{frame}{求解器与参数}
\begin{columns}[T]
\begin{column}{0.54\textwidth}
\textbf{外层 Krylov.}
\begin{itemize}\itemsep0.2em\footnotesize
  \item CG,\texttt{-ksp\_norm\_type preconditioned}
  \item \texttt{rtol}$=10^{-6}$,\texttt{atol}$=10^{-12}$,零初值
\end{itemize}
\vskip0.4em
\textbf{预条件子.}
\begin{itemize}\itemsep0.2em\footnotesize
  \item \texttt{-pc\_type asm -pc\_asm\_type basic}
  \item overlap $O\in\{0,1,2\}$
  \item 子域:\texttt{-sub\_pc\_type icc},填充 $L\in\{0,1,2\}$
\end{itemize}
\end{column}
\begin{column}{0.46\textwidth}
\textbf{两套独立实现.}
\begin{itemize}\itemsep0.2em\footnotesize
  \item MFEM 装配 $\to$ PETSc(\texttt{asm\_demo})
  \item 纯 PETSc P1 装配,不依赖 MFEM/HYPRE
  \item 用于排除"实现/装配产物"
\end{itemize}
\vskip0.4em
\textbf{扫描.} 对每个方案,扫 $O\times L$ 共 9 个组合,记录迭代次数、全局
归约次数(\texttt{-log\_view})与解的正确性。
\end{column}
\end{columns}
\end{frame}

% ---- 5. Background theory (condition number) ----
\begin{frame}{背景:单层加性 Schwarz 的条件数分解}
预条件 CG 的迭代次数由条件数支配,$n_{\text{iter}}\sim \tfrac12\sqrt{\kappa}\,\ln(2/\varepsilon)$,
$\kappa=\kappa(M^{-1}A)$。对单层加性 Schwarz,已有条件数上界(Toselli \& Widlund):
\[
\kappa(M^{-1}A)\ \lesssim\
\underbrace{C_0^{2}(\delta)}_{\substack{\text{稳定分解常数}\\ \text{随 overlap }\delta\,\text{增大而减小}}}
\ \cdot\
\underbrace{N_c}_{\substack{\text{子域着色数}\\ \text{(重叠区重复计数)}}}
\ \cdot\
\underbrace{\omega_{\max}/\omega_{\min}}_{\substack{\text{子域求解器谱等价}\\ \text{(精确解时}=1)}}
\]
\vskip0.6em
\begin{itemize}\itemsep0.3em
  \item 子域\textbf{精确}求解时,后两个因子为常数($\omega$ 比 $=1$),$C_0^2(\delta)$ 随 overlap 减小
        $\Rightarrow$ 迭代次数减少(经典结论)。
  \item 本工作考察:子域\textbf{不精确}(ICC)且采用 BASIC 加权时,后两个因子如何随 overlap 变化。
\end{itemize}
\end{frame}

% ---- 6. Methods: two single-level constructions ----
\begin{frame}{方法:两种单层构造}
保持外层为 CG、不引入两层(粗空间)方法。记 $R_i$ 为子域限制算子,$m_k$ 为自由度 $k$
所属的子域个数,$D=\operatorname{diag}(m_k)$。
\vskip0.6em
\textbf{(a) 对称缩放加性 Schwarz(sASM).}
\[
M^{-1}_{\text{sASM}}=D^{-1/2}\,M^{-1}_{\text{BASIC}}\,D^{-1/2},
\qquad M^{-1}_{\text{BASIC}}=\textstyle\sum_i R_i^{\top}\widetilde{A}_i^{-1}R_i .
\]
$D^{-1/2}$ 抵消重叠区的重复计数;$D$ 为对角阵,算子保持对称,故 CG 仍适用。
$O{=}0$ 时 $D=I$,退化为 BASIC。
\vskip0.6em
\textbf{(b) sASM $+$ Chebyshev 块解.}
将子域求解器从单次 ICC 改为固定 $k$ 步 Chebyshev 迭代(以 ICC 为光滑子,冻结特征值
区间)。这是对子域逆 $\widetilde{A}_i^{-1}$ 的多项式逼近,不增加填充/内存,且子域内运算
不含全局通信;固定步数的 Chebyshev 是固定的对称线性算子,CG 仍适用。
\end{frame}

% ---- 7. Result: subdomain solve accuracy (BASIC) ----
\begin{frame}{结果(一):子域求解精度的影响(固定 BASIC)}
\framesubtitle{主测问题,$n_x{=}48$,4 ranks;仅改变子域求解器,其余不变}
\begin{center}
\begin{tabular}{l ccc c}
\toprule
子域求解器 & $O=0$ & $O=1$ & $O=2$ & 随 overlap 趋势 \\
\midrule
精确 Cholesky & 52 & 34 & 30 & \da\ 下降 \\
ICC$(2)$ & 87 & 92 & 90 & 近持平 \\
ICC$(0)$ & 132 & 148 & 179 & \ua\ 上升 \\
\bottomrule
\end{tabular}
\end{center}
\vskip0.7em
\textbf{观察.} 子域精确求解时,overlap 增大使迭代次数下降(符合经典结论);
随子域求解精度降低(Cholesky $\to$ ICC$(2)$ $\to$ ICC$(0)$),趋势由下降转为上升。
即:overlap 增大导致迭代上升,需以\emph{子域求解不精确}为前提。
\end{frame}

% ---- 8. Result: three single-level preconditioners, full (O,L) ----
\begin{frame}{结果(二):三种单层预条件的完整结果}
\framesubtitle{主测问题,$n_x{=}48$,4 ranks;迭代次数,全部 $9$ 个 $(O,L)$ 组合}
\begin{center}
\begin{tabular}{cc ccc}
\toprule
$O$ & $L$ & scheme 0 & scheme 3 & scheme 4 \\
   &    & BASIC$+$ICC & sASM$+$ICC & sASM$+$Cheby \\
\midrule
0 & 0 & 132 & 132 & 85 \\
0 & 1 & 102 & 102 & 73 \\
0 & 2 & 87  & 87  & 66 \\
\midrule
1 & 0 & 148 & 104 & 68 \\
1 & 1 & 102 & 75  & 54 \\
1 & 2 & 92  & 70  & 47 \\
\midrule
2 & 0 & 179 & 103 & 65 \\
2 & 1 & 117 & 73  & 48 \\
2 & 2 & 90  & 67  & 42 \\
\bottomrule
\end{tabular}
\end{center}
\vskip0.4em
\footnotesize
\textbf{观察.} $L{=}0$ 列随 $O$ 的趋势:scheme 0 为 $132\!\to\!148\!\to\!179$(上升),
scheme 3 为 $132\!\to\!104\!\to\!103$(下降),scheme 4 为 $85\!\to\!68\!\to\!65$(下降)。
$O{=}0$ 时 scheme 0 与 scheme 3 相同($D{=}I$)。每个 $(O,L)$ 上 scheme 4 最低。
\end{frame}

% ---- 9. Verification: two implementations ----
\begin{frame}{验证:两套独立实现的一致性}
\framesubtitle{排除"实现/装配产物"}
\begin{columns}[T]
\begin{column}{0.5\textwidth}
\textbf{算子指纹}(同一确定性向量 $x$ 的 $\|Ax\|$):
\vskip0.3em
\begin{center}\footnotesize
\begin{tabular}{c c c}
\toprule
$n_x$ & MFEM $\|Ax\|$ & 纯 PETSc $\|Ax\|$ \\
\midrule
16 & 2.409188e$+$01 & 2.409188e$+$01 \\
17 & 2.137087e$+$01 & 2.137087e$+$01 \\
24 & 3.128297e$+$01 & 3.128297e$+$01 \\
\bottomrule
\end{tabular}
\end{center}
\end{column}
\begin{column}{0.5\textwidth}
\textbf{迭代次数一致性.}
\begin{itemize}\itemsep0.3em\footnotesize
  \item 三层对齐:矩阵零容差过滤、共享 METIS 分区、\texttt{MatLoad} 前预置逐进程行布局。
  \item 对齐后,scheme 0/3/4 在全部 $9$ 个 $(O,L)$ 上,两实现迭代次数逐位相同($27/27$)。
  \item \texttt{-ksp\_monitor} 下每一步预条件残差范数逐位相同。
\end{itemize}
\end{column}
\end{columns}
\vskip0.7em
\textbf{观察.} 该现象在不依赖 MFEM/HYPRE 的独立实现上同样出现,且两实现完全一致;
故为方法本身性质,非实现产物。
\end{frame}

% ---- 10. Universality: two other problems ----
\begin{frame}{结果(三):另两个问题上的普遍性}
\begin{columns}[T]
\begin{column}{0.52\textwidth}
\textbf{全 Neumann(奇异)椭圆}\ \footnotesize($n_x{=}24$,迭代次数)
\vskip0.2em
\begin{center}\footnotesize
\begin{tabular}{cc cc}
\toprule
$O$ & $L$ & BASIC & sASM \\
\midrule
0 & 0 & 64 & 64 \\
1 & 0 & 68 & 52 \\
2 & 0 & 87 & 51 \\
\midrule
0 & 2 & 50 & 50 \\
1 & 2 & 45 & 37 \\
2 & 2 & 44 & 34 \\
\bottomrule
\end{tabular}
\end{center}
\end{column}
\begin{column}{0.48\textwidth}
\textbf{reaction--diffusion}\ \footnotesize($\tfrac1{\Delta t}M{+}\tfrac12 K$,$n_x{=}48$,BASIC,迭代次数)
\vskip0.2em
\begin{center}\footnotesize
\begin{tabular}{c ccc}
\toprule
 & $O=0$ & $O=1$ & $O=2$ \\
\midrule
ICC$(0)$,\,$\Delta t{=}10^{-2}$ & 33 & 42 & 47 \\
ICC$(0)$,\,$\Delta t{=}10^{-3}$ & 12 & 15 & 18 \\
ICC$(0)$,\,$\Delta t{=}10^{-4}$ & 6  & 8  & 9  \\
\midrule
ICC$(2)$,\,$\Delta t{=}10^{-4}$ & 6  & 6  & 7  \\
\bottomrule
\end{tabular}
\end{center}
\end{column}
\end{columns}
\vskip0.7em
\textbf{观察.} 全 Neumann 问题与质量主导的 reaction--diffusion 问题,在子域求解不精确
(低 ICC)时同样出现 overlap 增大、迭代上升;sASM 与提高 ICC 填充均使其转为下降。
reaction--diffusion 因质量项良态,绝对迭代次数较小。
\end{frame}

% ---- 11. Synchronisation & time ----
\begin{frame}{结果(四):全局同步次数与墙钟时间}
\framesubtitle{主测问题,$n_x{=}48$,4 ranks}
\begin{columns}[T]
\begin{column}{0.5\textwidth}
\textbf{全局归约次数}\ \footnotesize(\texttt{-log\_view},$L{=}0$)
\vskip0.2em
\begin{center}\footnotesize
\begin{tabular}{l ccc}
\toprule
方案 & $O{=}0$ & $O{=}1$ & $O{=}2$ \\
\midrule
scheme 0 & 525 & 574 & 668 \\
scheme 3 & 528 & 445 & 443 \\
scheme 4 & 387 & 337 & 329 \\
\bottomrule
\end{tabular}
\end{center}
\footnotesize 每次外层迭代含 2 次点积归约;内层 Chebyshev 在子域(\texttt{COMM\_SELF})上,不产生全局归约。
\end{column}
\begin{column}{0.5\textwidth}
\textbf{墙钟时间}\ \footnotesize($O{=}1$,3 次中位数,秒)
\vskip0.2em
\begin{center}\footnotesize
\begin{tabular}{l ccc}
\toprule
方案 & $L{=}0$ & $L{=}1$ & $L{=}2$ \\
\midrule
scheme 0 & 0.231 & 0.168 & 0.171 \\
scheme 3 & 0.149 & 0.112 & 0.119 \\
scheme 4 & 0.130 & 0.112 & 0.113 \\
\bottomrule
\end{tabular}
\end{center}
\footnotesize 墙钟时间在负载机器上不稳定,仅供参考;全局归约次数为确定性指标。
\end{column}
\end{columns}
\vskip0.6em
\textbf{观察.} 在 $O{=}2$,$L{=}0$ 处,全局归约次数由 $668$(scheme 0)降至 $329$(scheme 4)。
\end{frame}

% ---- 12. Relation to literature & theory ----
\begin{frame}{与文献和理论的关系}
\begin{columns}[T]
\begin{column}{0.5\textwidth}
\textbf{现象的相关文献.}
\begin{itemize}\itemsep0.25em\footnotesize
  \item Cai \& Sarkis (1999):RAS,指出 BASIC 在重叠区的重复计数。
  \item Efstathiou \& Gander (2003):加性 Schwarz 的收敛依赖子域着色数。
  \item Ernst, Flemisch \& Wohlmuth (2009):记录了"overlap 增大时不精确方法迭代次数较精确方法更差"(非线性 multiplicative 设定)。
\end{itemize}
\end{column}
\begin{column}{0.5\textwidth}
\textbf{理论状态.}
\vskip0.2em
\begin{center}\footnotesize
\begin{tabular}{>{\raggedright\arraybackslash}p{3.3cm} >{\raggedright\arraybackslash}p{2.5cm}}
\toprule
命题 & 状态 \\
\midrule
$\lambda_{\max}\le N_c$(重复计数) & 已证 \\
不精确解的 $\omega$ 谱等价因子 & 已证 \\
精确解时 overlap 增大$\Rightarrow$迭代降 & 已证 \\
缩放使 $\lambda_{\max}\le 1$(sASM) & 已证 \\
不精确 $\times$ 重复计数 $\Rightarrow$ 随 overlap 严格上升 & 机制与实证,尚无闭式定理 \\
\bottomrule
\end{tabular}
\end{center}
\end{column}
\end{columns}
\vskip0.6em
\footnotesize
未决问题:给出固定填充 ICC 在增大子块上的谱等价常数 $\omega(\delta)$ 随 overlap 增长的定量(闭式)下界。
\end{frame}

% ---- 13. Summary ----
\begin{frame}{小结}
\begin{itemize}\itemsep0.5em
  \item 在 BASIC$+$ICC$+$CG 下,椭圆问题的 CG 迭代次数随 overlap 增大而增加;该现象在
        两套独立实现上逐位一致,为方法本身性质。
  \item 该现象需要两个条件同时成立:子域求解不精确,以及 BASIC 在重叠区的重复计数。
        去除任一(精确子域解,或 sASM 缩放)即恢复 overlap 增大、迭代减少的经典趋势。
  \item sASM 与 sASM$+$Chebyshev 均保持外层 CG 对称、保持单层;在主测问题上同时减少
        迭代次数与全局归约次数。
  \item 同一行为在三个模型问题上复现(混合边界 Laplace、全 Neumann 奇异、reaction--diffusion)。
  \item 两个因素已分别有理论刻画;两者共同导致"随 overlap 严格上升"的闭式证明仍未决。
\end{itemize}
\end{frame}

\end{document}
